
\chapter*{Acknowledgments}
Águeda Gómez Cambronero\\
Silvia Ferrer Castillo y Marta Mestre Aguila\\
Fabio Nicolas Luna Arevalo

\clearpage

\chapter*{Abstract}
This document presents the Bachelor’s Thesis developed for the Bachelor’s Degree in Video Game Design and Development. The project originates from a pedagogical activity carried out with Early Childhood Education students in the course MI1815 – The Natural Environment in Early Childhood Education, consisting of a guided visit to the \textit{Jardín de los Sentidos}, an outdoor space at Jaume I University characterised by its rich biodiversity. The educational value of this activity lies in encouraging students to observe, discuss, and reflect on natural ecosystems through direct interaction with the environment.
Building upon this experience, this project explores how educational objectives of the original activity can be transferred to a digital, autonomous and playful format. To address this challenge, \textit{Micorrizas} was designed as an educational video game that promotes environmental awareness and engagement through active exploration and collaborative learning.
The resulting game is a cooperative location-based mobile experience developed with Unity, where players interact with real-world locations while progressing through shared challenges and decision-making processes. To support this experience, the game combines geolocation and multiplayer technologies that enable collaborative gameplay in outdoor environments. The game has also been evaluated by the lecturers responsible for the original educational activity, providing feedback on its suitability as a tool for environmental education in outdoor learning contexts.

\vspace{1cm}
\section*{Keywords}
Serious game - educational game - cooperative multiplayer - geolocation - mobile learning - environmental education - Jardín de los Sentidos.

\clearpage

\tableofcontents
\clearpage

\listoffigures
\clearpage

\listoftables
\clearpage

